% Exam Template for UMTYMP and Math Department courses
%
% Using Philip Hirschhorn's exam.cls: http://www-math.mit.edu/~psh/#ExamCls
%
% run pdflatex on a finished exam at least three times to do the grading table on front page.
%
%%%%%%%%%%%%%%%%%%%%%%%%%%%%%%%%%%%%%%%%%%%%%%%%%%%%%%%%%%%%%%%%%%%%%%%%%%%%%%%%%%%%%%%%%%%%%%

% These lines can probably stay unchanged, although you can remove the last
% two packages if you're not making pictures with tikz.
\documentclass[11pt]{exam}
\RequirePackage{amssymb, amsfonts, amsmath, latexsym, verbatim, xspace, setspace}
\RequirePackage{tikz, pgflibraryplotmarks}

% By default LaTeX uses large margins.  This doesn't work well on exams; problems
% end up in the "middle" of the page, reducing the amount of space for students
% to work on them.
\usepackage[margin=1in]{geometry}
\usepackage{enumerate}
\usepackage{amsthm}

\theoremstyle{definition}
\newtheorem{soln}{Solution}

% Here's where you edit the Class, Exam, Date, etc.
\newcommand{\class}{Math 350 Section 1}
\newcommand{\term}{Fall 2024}
\newcommand{\examnum}{Exam II}
\newcommand{\examdate}{October 23, 2024}
\newcommand{\timelimit}{1 Hour 50 Minutes}
\newcommand{\ol}[1]{\overline{#1}}

% For an exam, single spacing is most appropriate
\singlespacing
% \onehalfspacing
% \doublespacing

% For an exam, we generally want to turn off paragraph indentation
\parindent 0ex

\begin{document} 

% These commands set up the running header on the top of the exam pages
\pagestyle{head}
\firstpageheader{}{}{}
\runningheader{\class}{\examnum\ - Page \thepage\ of \numpages}{\examdate}
\runningheadrule

\begin{flushright}
\begin{tabular}{p{2.8in} r l}
\textbf{\class} & \textbf{Name (Print):} & \makebox[2in]{\hrulefill}\\
\textbf{\term} &&\\
\textbf{\examnum} & \textbf{Student ID:}&\makebox[2in]{\hrulefill}\\
\textbf{\examdate} &&\\
\textbf{Time Limit: \timelimit} % & Teaching Assistant & \makebox[2in]{\hrulefill}
\end{tabular}\\
\end{flushright}
\rule[1ex]{\textwidth}{.1pt}


This exam contains \numpages\ pages (including this cover page) and
\numquestions\ problems.  Check to see if any pages are missing.  Enter
all requested information on the top of this page, and put your initials
on the top of every page, in case the pages become separated.\\

You may \textit{not} use your books or notes on this exam.

You are required to show your work on each problem on this exam.  The following rules apply:\\

\begin{minipage}[t]{3.7in}
\vspace{0pt}
\begin{itemize}

%\item \textbf{If you use a ``fundamental theorem'' you must indicate this} and explain
%why the theorem may be applied.

\item \textbf{Organize your work}, in a reasonably neat and coherent way, in
the space provided. Work scattered all over the page without a clear ordering will 
receive very little credit.  

\item \textbf{Mysterious or unsupported answers will not receive full
credit}.  A correct answer, unsupported by calculations, explanation,
or algebraic work will receive no credit; an incorrect answer supported
by substantially correct calculations and explanations might still receive
partial credit.  This especially applies to limit calculations.

\item If you need more space, use the back of the pages; clearly indicate when you have done this.

\item \textbf{Box Your Answer} where appropriate, in order to clearly indicate what you consider the answer to the question to be.
\end{itemize}

Do not write in the table to the right.
\end{minipage}
\hfill
\begin{minipage}[t]{2.3in}
\vspace{0pt}
%\cellwidth{3em}
\gradetablestretch{2}
\vqword{Problem}
\addpoints % required here by exam.cls, even though questions haven't started yet.	
\gradetable[v]%[pages]  % Use [pages] to have grading table by page instead of question

\end{minipage}
\newpage % End of cover page

%%%%%%%%%%%%%%%%%%%%%%%%%%%%%%%%%%%%%%%%%%%%%%%%%%%%%%%%%%%%%%%%%%%%%%%%%%%%%%%%%%%%%
%
% See http://www-math.mit.edu/~psh/#ExamCls for full documentation, but the questions
% below give an idea of how to write questions [with parts] and have the points
% tracked automatically on the cover page.
%

%%%%%%%%%%%%%%%%%%%%%%%%%%%%%%%%%%%%%%%%%%%%%%%%%%%%%%%%%%%%%%%%%%%%%%%%%%%%%%%%%%%%%

\begin{questions}

\addpoints

\question[10]\mbox{}
\textbf{TRUE or FALSE!}  Write  TRUE if the statement is true.  Otherwise, write FALSE.  Your response should be in ALL CAPS.  No justification is required.
\begin{enumerate}[(a)]
\item a monotone increasing sequence must converge 
\vspace{1.2in}
\item the union of any family of open sets is open
\vspace{1.2in}
\item if a set is not closed, then it must be open
\vspace{1.2in}
\item in a metric space with the discrete metric, every set is both open and closed
\vspace{1.2in}
\item every convergent sequence is also Cauchy
\vspace{1.2in}
\end{enumerate}

\newpage
\question[10]\mbox{}

\begin{enumerate}[(a)]
\item
Write down the definition of a metric space.
\vspace{3in}
\item 
Write down three different examples of metrics we can define on $\mathbb{R}^2$.
\vspace{2in}
\item 
Explain why $d(x,y) = |x+y|$ is NOT a metric on $\mathbb{R}$.
\end{enumerate}

\newpage
\question[10]\mbox{}

\begin{enumerate}[(a)]
\item
Write down the definition of an open set in a metric space $(M,d)$.
\vspace{1in}
\item 
Let $(M,d)$ be a metric space and let $x\in M$ and $r>0$.
Prove that the open ball
$$B_M(x;r)=\{y\in M: d(x,y) < r\}$$
is an open set without using any theorems from class.
\end{enumerate}

\newpage
\question[10]\mbox{}
\begin{enumerate}[(a)]
\item Write down the definition of a sequence $\{x_n\}$ of values in a metric space $(M,d)$ converging to a value $L\in M$.
\vspace{1in}

\item Prove that if a sequence $\{x_n\}$ converges to a value $L\in M$, then there exists a constant $R>0$ such that $d(x_n,L) < R$ for all $n\in\mathbb{Z}_+$.
\end{enumerate}

\newpage
\question[10]\mbox{}
Carefully state each of the following theorems.
\begin{enumerate}[(a)]
\item the Lindel\"{o}f Covering Theorem
\vspace{2in}
\item the Bolzano-Weierstrass Theorem
\vspace{2in}
\item the Cantor Intersection Theorem
\vspace{2in}
\item the Heine-Borel Theorem
\end{enumerate}

\newpage
\question[10]\mbox{}
\begin{enumerate}[(a)]
\item Write down the definition of an open cover of a subset $A\subseteq M$ of a metric space $M$.
\vspace{2in}
\item Prove that the family of sets $\{U_i: i\in I\}$ where $I=\mathbb{Z}_+$ and $U_i = (1/i,2)$ is an open cover of the interval $(0,1]$
\end{enumerate}

\newpage
\question[10]\mbox{}
\begin{enumerate}[(a)]
\item Write down the definition of a subset $A\subseteq M$ of a metric space $M$ being compact.
\vspace{2in}
\item Prove that if $A = \{a\}$ for some $a\in M$, then $A$ is compact.
\end{enumerate}


\end{questions}

\end{document}

